% =====================================================================
%  APPENDIX: The non-Archimedean shadow -- supertasks, closed causal
%  loops, and the division-algebra ledger at the finite places
%  ADDED 2026-07-23. Writes up, as an invitation to the community, the
%  p-adic / CTC / division-algebra correspondences developed in the
%  session of 2026-07-23. Companion to the open-problem document
%  open_problem_adelic_mass_constraints.md (whose Problems 2, 4(b),
%  4(c) and minimal-progress criterion (B) are engaged below, without
%  upgrading the status of its conjectures) and to the continuity
%  appendix app:continuity-open, whose Accumulation Lemma
%  (lem:accumulation), separating witnesses (rem:witnesses),
%  Archimedean fork (app:cont-scalar) and dyadic-tower construction
%  (eq:dyadic-tower) are cross-referenced and MUST be present in the
%  same document.
%  Requires (amsthm): lemma, theorem, corollary, remark, openproblem
%  as declared for app:continuity-open, plus
%  \newtheorem{proposition}[theorem]{Proposition} if not already
%  declared (VERIFY against the master preamble).
%  External labels assumed from the master file: sec:self-reference,
%  sec:fanout-boundary, sec:born, sec:hopf (VERIFY the last two --
%  the same mapping task as the Jaynes appendix).
%  Citations use \autocite per house convention. Keys tagged NEW in
%  the trailer block are supplied in the accompanying batch
%  refs_padic_ctc_division_2026-07-23T0617.ris; keys tagged EXISTING
%  are assumed already in the master bibliography (VERIFY each).
% =====================================================================
%
% ---------------------------------------------------------------------
%  MAIN-BODY INSERT (place in sec:self-reference, at the discussion of
%  negative probability and supertasks; adjust the lead-in):
%
%  The analytic continuation that carries Shannon's measure to complex
%  and negative values is the Archimedean route to a pre-record
%  ledger. Ostrowski's fork offers a second route: Khrennikov's
%  p-adic probability theory realises negative probabilities as
%  honest limits of relative frequencies in the non-Archimedean
%  completions, where a negative integer is written with infinitely
%  many digits to the left --- the completed value of a supertask of
%  borrowing. Appendix~\ref{app:padic-shadow} develops the
%  correspondence: which supertasks converge at which places, why the
%  closed-loop reading of measurement acquires contractive fixed
%  points p-adically, how Hilbert reciprocity forces the
%  division-algebra obstructions of the main text to balance across
%  the places with the prime two as Hamilton's partner, and six open
%  problems inviting the p-adic and division-algebra communities to
%  meet.
% ---------------------------------------------------------------------

\section{The non-Archimedean shadow: supertasks, closed loops, and the division-algebra ledger at the finite places}
\label{app:padic-shadow}

The framework's record layer is Archimedean by argument, not habit:
Lemma~\ref{lem:accumulation} selects the real completion of $\mathbb{Q}$ because the
Thomson halving schedule accumulates there and nowhere else, and the exclusions of
Remark~\ref{rem:witnesses} then consume those real scalars. But the road not taken at
Ostrowski's fork is not empty. Khrennikov's $p$-adic probability theory
\autocite{khrennikov1995} is its consistent mathematics, realising negative probabilities
as honest limits of relative frequencies in the non-Archimedean topologies --- with the
striking feature, noted in the main text, that a negative integer's $p$-adic expansion
carries infinitely many digits to the \emph{left}, so that writing down $-1$ is already
the completed value of a supertask of borrowing. Historically the two tines were offered
as rivals, $p$-adic quantum mechanics competing with complex quantum mechanics for the
role of true home of the theory, and the proposal found little uptake. This appendix
takes the opposite view: the tines have complementary jobs. Records, whose statistics
are counts of copies, must live at the Archimedean place; the pre-record ledger, where
the framework's sign cancellations are thermodynamically free, owes Archimedes nothing;
and the adelic product formula $\prod_v |x|_v = 1$ \autocite{ostrowski1916} is the
consistency condition binding the two descriptions of one bookkeeping, in the register
of the adelic programme of Vladimirov, Volovich and Dragovich
\autocite{vladimirovvolovich1989, dragovich1995}.

What follows contains two elementary lemmas and two propositions (proved), one theorem
compiled from classical local--global algebra (cited, not reproved), one corollary, and
six open problems. As throughout the corpus, RHYME marks a resemblance offered for
interrogation, not evidence; CANDIDATE marks a framework reading of a theorem; nothing
below upgrades the status of the adelic mass-consistency conjecture, though one of its
minimal-progress criteria is discharged by citation in \S\ref{app:padic-ledger}. The
problems are stated so that a specialist can engage each without adopting the
framework's other commitments, in the manner of
Appendix~\ref{app:continuity-open}.

\subsection{Two supertasks, two tines}
\label{app:padic-supertasks}

Lemma~\ref{lem:accumulation} concerns the halving schedule: infinitely many digits
completed to the \emph{right}, convergent at the Archimedean place alone. Its dual
concerns doubling.

\begin{lemma}[The doubling schedule converges only $p$-adically]\label{lem:doubling}
Fix a prime $p$ and consider $\sum_{k \ge 0} (p-1)\,p^{k}$, with partial sums
$s_n = p^{n+1} - 1$. The series diverges in $\mathbb{R}$ and in every
$\mathbb{Q}_{\ell}$ with $\ell \ne p$; in $\mathbb{Q}_p$ it converges, with
\begin{equation}
\sum_{k=0}^{\infty} (p-1)\,p^{k} \;=\; -1,
\qquad\text{in particular}\qquad
1 + 2 + 4 + 8 + \cdots \;=\; -1 \ \text{ in } \mathbb{Q}_2.
\label{eq:fanout-tally}
\end{equation}
Moreover $p^{n} \to 0$ in $\mathbb{Q}_p$.
\end{lemma}

\begin{proof}
Real divergence: the partial sums are monotone and unbounded. In $\mathbb{Q}_{\ell}$
with $\ell \ne p$ the terms have constant nonzero absolute value
$|(p-1)p^{k}|_{\ell} = |p-1|_{\ell}$, and a series in a non-Archimedean field converges
only if its terms tend to zero \autocite{gouvea1997}. In $\mathbb{Q}_p$,
$|s_n - (-1)|_p = |p^{n+1}|_p = p^{-(n+1)} \to 0$.
\end{proof}

\begin{remark}[Digits and directions]\label{rem:digits}
A real number completes infinitely many digits rightward; the Thomson schedule of
Lemma~\ref{lem:accumulation} is that completion read as a schedule
\autocite{thomson1954}. A $p$-adic integer completes infinitely many digits leftward;
the expansion $-1 = \cdots(p-1)(p-1)(p-1)$ is Lemma~\ref{lem:doubling} read as digits,
the carry propagated forever. Ostrowski's fork is thus a choice of \emph{direction}:
which of the two digit-infinities is permitted to terminate. A negative integer,
$p$-adically, is not a primitive but the completed value of a supertask of borrowing
--- creation from nothing, in the sense of the classical supertask literature
\autocite{laraudogoitia1998}, in arithmetic rather than mechanical dress.
\end{remark}

\begin{lemma}[FANOUT-blindness of the finite places]\label{lem:fanout-blind}
An absolute value on $\mathbb{Q}$ is non-Archimedean if and only if $|n| \le 1$ for
every integer $n$.
\end{lemma}

\begin{proof}
Standard; the nontrivial direction is the binomial estimate \autocite{gouvea1997}.
\end{proof}

The Archimedean axiom is the statement that repetition accumulates:
$|1 + 1 + \cdots + 1|$ is unbounded. Repetition of a token is precisely what FANOUT
(\S\ref{sec:fanout-boundary}) outputs, and frequencies are counts of copies. The record
layer is therefore not merely permitted to be Archimedean; counting forces it --- the
Accumulation Lemma seen from the copy side. Conversely, by Lemma~\ref{lem:fanout-blind}
the finite places are \emph{FANOUT-blind}: copying never increases magnitude there, and
by Lemma~\ref{lem:doubling} an unbounded doubling cascade has, $2$-adically, a record
count tending to zero and a cumulative tally converging to minus one.
Equation~\eqref{eq:fanout-tally} is this appendix's emblem: unbounded copying, tallied
at the FANOUT-blind place proper to doubling, is finite, convergent --- and negative.
We record the identification as STRUCTURAL: it rereads Ostrowski's dichotomy,
Khrennikov's frequency theory (collectives whose relative frequencies stabilise
$p$-adically though not really \autocite{khrennikov1995}), and the framework's Landauer
accounting (extensive erasure costs are Archimedean counting; pre-FANOUT sign
cancellation is free exactly where repetition does not accumulate) as three faces of
one fork.

\begin{openproblem}[Measurement as Archimedeanisation]\label{op:padic-frequency}
Formulate the FANOUT criterion inside $p$-adic frequency theory: show that a collective
admits Archimedean stabilisation of its relative frequencies precisely when its
outcomes are copied records, and exhibit at least one of the framework's pre-FANOUT
sign cancellations as a collective that stabilises $p$-adically in Khrennikov's sense.
Progress in either direction is informative. A worked example would give the
pre-record ledger a non-Archimedean probability semantics requiring no analytic
continuation; a proof that the framework's cancellations never stabilise $p$-adically
would separate the two routes to negative probability cleanly.
\end{openproblem}

\subsection{One fixed point, two temperaments}
\label{app:padic-carry}

\begin{proposition}[The carry map]\label{prop:carry}
Let $f(x) = px + (p-1)$. On $\mathbb{Z}_p$, $f$ is a contraction with constant
$p^{-1}$; its unique fixed point is $-1$, and the orbit of $0$ is
$f^{n}(0) = p^{n} - 1 \to -1$. On $\mathbb{R}$, the same map is expanding with constant
$p$; the fixed point $-1$ is repelling, and the orbit of $0$ diverges.
\end{proposition}

\begin{proof}
$|f(x) - f(y)|_p = |p|_p\,|x - y|_p = p^{-1}|x - y|_p$; Banach's theorem on the
complete ultrametric space $\mathbb{Z}_p$ gives a unique fixed point, which solves
$x = px + (p-1)$, that is $x = -1$; induction gives $f^{n}(0) = p^{n} - 1$. Over
$\mathbb{R}$ the Lipschitz constant is $p > 1$ and the orbit of $0$ is
$p^{n} - 1 \to \infty$.
\end{proof}

\begin{remark}[Hensel]\label{rem:hensel}
Hensel's lemma is the general machine of the same kind: approximate self-consistency
modulo $p$ lifts, by ultrametric contraction, to exact self-consistency in
$\mathbb{Z}_p$ \autocite{gouvea1997, serre1973}.
\end{remark}

The main text models a measurement outcome as computationally indistinguishable from
information generated at the fixed point of a closed causal loop
(\S\ref{sec:self-reference}; \autocite{deutsch1991, aaronsonwatrous2009}). Deutsch's
consistency condition $\rho = \Phi(\rho)$ is solved over $\mathbb{R}$ by a compactness
argument --- Brouwer's theorem on the convex body of density operators --- which
certifies existence without exhibiting convergence; the supertask reading of the same
situation \autocite{laraudogoitia1998} is exactly the divergent Archimedean iteration.
Proposition~\ref{prop:carry} displays the two temperaments on the simplest
self-consistency equation there is: one equation, one fixed point, repelling at the
Archimedean place and attracting at the finite one. We fence the reading as RHYME:
\emph{what is a paradox at the Archimedean place is an attractor at the finite
places.} A closed timelike curve and a $p$-adic metric are, on this reading, two
devices for rendering a self-referential equation consistent --- the first geometric
(impose the fixed point on a loop), the second analytic (complete in the topology
where the iteration contracts).

\begin{openproblem}[Deutsch consistency as Hensel lifting]\label{op:padic-hensel}
Characterise the Deutsch-type self-consistency maps $\Phi$ that become contractions in
a suitable $p$-adic operator framework, so that their fixed points are delivered
constructively by Banach--Hensel iteration rather than existentially by Brouwer. The
$p$-adic dynamics literature \autocite{khrennikovnilsson2004} classifies attracting
fixed points of polynomial and monomial maps on $\mathbb{Z}_p$; the question is which
physically meaningful consistency conditions fall in the attracting class, with
Proposition~\ref{prop:carry} as the seed example. A positive answer would give the
closed-loop reading of measurement a constructive backend at the finite places.
\end{openproblem}

\subsection{The division-algebra ledger}
\label{app:padic-ledger}

The main text's spine runs from closed-loop consistency through parallelisable spheres
to the Hurwitz tower $\mathbb{R} \subset \mathbb{C} \subset \mathbb{H} \subset
\mathbb{O}$ \autocite{hurwitz1898, adams1960}. That tower has always been read at the
real place. Classical local--global algebra reads it at all places at once, and the
result is a ledger. Write $(a,b)_v$ for the Hilbert symbol of
$a, b \in \mathbb{Q}^{\times}$ at a place $v$.

\begin{theorem}[The ledger at the finite places; classical]\label{thm:ledger}
\begin{itemize}
\item[\textup{(i)}] \emph{Reciprocity.} $\prod_v (a,b)_v = 1$: the product formula for
Hilbert symbols, equivalent to quadratic reciprocity \autocite{serre1973}.
\item[\textup{(ii)}] \emph{Hamilton's partner.} $(-1,-1)_v = -1$ exactly for
$v \in \{2, \infty\}$; hence the rational Hamilton algebra $(-1,-1)_{\mathbb{Q}}$ is a
division algebra at precisely the real place and the prime two \autocite{serre1973}.
\item[\textup{(iii)}] \emph{Even ramification.} A quaternion algebra over $\mathbb{Q}$
is a division algebra at a finite set of places of even cardinality, and is determined
up to isomorphism by that set (Albert--Brauer--Hasse--Noether)
\autocite{brauerhassenoether1932, alberthasse1932}.
\item[\textup{(iv)}] \emph{Local uniqueness.} Each $\mathbb{Q}_p$ carries exactly one
quaternion division algebra up to isomorphism; for odd $p$ it may be presented as
$(u, p)$ with $u$ a non-square unit \autocite{serre1973}.
\item[\textup{(v)}] \emph{No local octonions.} Every quadratic form over $\mathbb{Q}_p$
in five or more variables is isotropic \autocite{serre1973}; octonion algebras over a
field are classified by their norm forms and are division precisely when the norm form
is anisotropic \autocite{jacobson1958, springerveldkamp2000}; hence every octonion
algebra over every $\mathbb{Q}_p$ is split.
\item[\textup{(vi)}] \emph{Archimedean exclusivity.} Over $\mathbb{Q}$ there are
exactly two octonion algebras --- the split algebra and the classical division
octonions --- and they are distinguished at the real place alone
\autocite{springerveldkamp2000}.
\end{itemize}
\end{theorem}

\begin{corollary}[The standard tower at the finite places]\label{cor:cd-tower}
Run the Cayley--Dickson construction over $\mathbb{Q}_p$ with parameter $-1$ at every
doubling. For odd $p$ the tower splits already at the quaternion rung, since
$(-1,-1)_p = 1$; for $p = 2$ it is a division tower through the quaternions and splits
at the octonion rung. At no finite place does the tower reach a division octonion
algebra.
\end{corollary}

\begin{proof}
Immediate from Theorem~\ref{thm:ledger}(ii) and (v).
\end{proof}

Three readings. \emph{First, a criterion discharged.}
Corollary~\ref{cor:cd-tower} settles, by citation rather than new computation, the
minimal-progress criterion of the adelic constraints programme that asked at which
stage the Cayley--Dickson construction over $\mathbb{Q}_2$ ceases to produce division
algebras --- and it shows the failure mode differs in kind from the real case. At the
real place the tower's termination is topological, Hopf invariant one
\autocite{adams1960}, and strikes one rung \emph{later}, at the sedenions; at the
finite places it is arithmetic, isotropy forced by the $u$-invariant
$u(\mathbb{Q}_p) = 4$ \autocite{lam2005, serre1973}, and strikes one rung
\emph{earlier}. Whether the Hopf-invariant-one obstruction has a motivic or etale
avatar at the finite places we leave untouched.

\emph{Second, conservation of obstruction.} Reciprocity, Theorem~\ref{thm:ledger}(i),
is a conservation law: the local obstructions of a quaternion algebra multiply to one,
so obstruction is never created or destroyed, only relocated among the places --- the
corpus's bump in the carpet, adelic edition. Hamilton's noncommutativity debt at the
real place is necessarily paired, and its partner, by (ii), is the prime two: the
binary place. The octonionic obstruction, by contrast, cannot be shared: (v) and (vi)
locate it at the Archimedean place alone.

\emph{Third, a CANDIDATE reading in the dictionary.} In the entanglement dictionary of
the main text the quaternionic layer carries the bipartite and measurement structure
($S^3$, $SU(2)$) and the octonionic layer the irreducibly tripartite GHZ structure.
The ledger then says: the quaternionic layer possesses a $2$-adic shadow --- the
compact, totally disconnected unit group of the $2$-adic quaternion division algebra
--- while the GHZ layer possesses none. CANDIDATE, since what is being read is the
dictionary entry, not the algebra.

\begin{openproblem}[A $p$-adic asymmetry between CHSH and GHZ]\label{op:padic-ghz}
Determine whether irreducibly tripartite (GHZ-class) correlations admit $p$-adic
probabilistic models in the frequency sense, as bipartite interference does in
Khrennikov's analyses \autocite{khrennikov2009}. The ledger predicts an asymmetry:
obstruction-sharing with a finite place is available to the quaternionic layer and
unavailable to the octonionic one. Either outcome is progress. An obstruction theorem
would be the first physical trace of the octonions' non-Archimedean splitting; a model
would refute the CANDIDATE reading above, and is the honest falsifier.
\end{openproblem}

\subsection{A $p$-adic Hopf quotient; positivity as the exotic phenomenon}
\label{app:padic-hopf}

\begin{proposition}[A Hopf-type quotient at $p$]\label{prop:padic-hopf}
Let $S = \mathbb{Z}_p^2 \setminus (p\mathbb{Z}_p)^2$ be the set of unimodular pairs.
The map $S \to \mathbb{P}^1(\mathbb{Q}_p)$, $(x,y) \mapsto [x:y]$, is a surjection
whose fibres are exactly the orbits of the unit group $\mathbb{Z}_p^{\times}$ acting
by scaling, so $S$ is a principal $\mathbb{Z}_p^{\times}$-bundle over
$\mathbb{P}^1(\mathbb{Q}_p)$. This bundle is trivial: the base is compact and
zero-dimensional, and admits a clopen partition over which unimodular sections exist.
Finally, the complete invariant of the $\mathbb{Z}_p^{\times}$-scaling action on
$\mathbb{Q}_p$ is the absolute value: $x$ and $y$ lie in one orbit if and only if
$|x|_p = |y|_p$.
\end{proposition}

\begin{proof}
Every point of $\mathbb{P}^1(\mathbb{Q}_p)$ has a representative with
$\max(|x|_p, |y|_p) = 1$, unique up to $\mathbb{Z}_p^{\times}$. The sets
$\{|x|_p \le |y|_p\}$ and $\{|x|_p > |y|_p\}$ are clopen and carry the sections
$[x:y] \mapsto (x/y,\, 1)$ and $[x:y] \mapsto (1,\, y/x)$ respectively. The orbits of
$\mathbb{Z}_p^{\times}$ on $\mathbb{Q}_p^{\times}$ are the spheres
$\{\,|x|_p = p^{-k}\,\}$, $k \in \mathbb{Z}$, together with the fixed point $0$.
\end{proof}

At the real place the Born derivation of the main text (\S\ref{sec:born}) quotients
the $U(1)$ ignorance fibre of the Hopf bundle (\S\ref{sec:hopf}); the surviving
invariant is quadratic, and positive. Proposition~\ref{prop:padic-hopf} runs the same
quotient at a finite place: the phase group is $\mathbb{Z}_p^{\times}$ --- for odd $p$
a torsion clock times a pro-$p$ depth,
$\mathbb{Z}_p^{\times} \cong \mu_{p-1} \times (1 + p\mathbb{Z}_p)$
\autocite{gouvea1997} --- and the Born datum it leaves is the valuation $|x|_p$. On
rational data the product of all Born data, real and $p$-adic, is one: the product
formula as a distributed Born ledger, in the register of adelic quantum mechanics
\autocite{vladimirovvolovich1989, dragovich1995}. Two inversions follow.

\emph{Positivity.} By Artin--Schreier a field is orderable if and only if it is
formally real \autocite{artinschreier1927}; $\mathbb{R}$ is the unique orderable
completion of $\mathbb{Q}$, while every $\mathbb{Q}_p$ has finite level --- minus one
is a sum of at most four squares there \autocite{lam2005} --- so squared magnitudes
carry no sign at any finite place. Non-negativity of probability is a theorem
available at exactly one place of $\mathbb{Q}$. From the non-Archimedean side it is
positivity, not negativity, that is the exotic phenomenon; negativity is
contextuality's Archimedean face \autocite{spekkens2008}, and the finite places never
had the positivity to lose.

\emph{The twist.} The bundle of Proposition~\ref{prop:padic-hopf} is trivial: a
zero-dimensional base cannot support the obstruction that, at the real place, encodes
the irrecoverability of the self-referential phase. The Hopf twist has no local avatar
at any finite place. Its conservation law is global: nontriviality at the real place
is the statement $(-1,-1)_{\infty} = -1$, and reciprocity,
Theorem~\ref{thm:ledger}(i)--(ii), forces the balancing entry at $p = 2$. We fence the
reading as RHYME: \emph{the Hopf twist obeys Gauss reciprocity.} The prime two will
appear once more, by duality, in \S\ref{app:padic-solenoid}.

\begin{openproblem}[The twist's conservation law]\label{op:padic-reciprocity}
Make the rhyme a theorem or refute it: identify the nontriviality class of the
Archimedean Hopf bundle with the infinite component of a global reciprocity datum ---
the class of $(-1,-1)_{\mathbb{Q}}$ in the $2$-torsion of the Brauer group is the
obvious candidate --- so that the $2$-adic component is the forced complement of the
observer's phase ignorance, and the triviality of Proposition~\ref{prop:padic-hopf} at
every finite place is the local expression of a globally conserved twist. Refutation
would consist in exhibiting a precise sense in which the bundle-theoretic and
Brauer-theoretic obstructions come apart.
\end{openproblem}

\subsection{The Tsirelson bound as ramification at two}
\label{app:padic-tsirelson}

For $\pm 1$-valued observables $A, A'$ and $B, B'$ on two sites, the CHSH operator
$\mathcal{C} = A \otimes B + A \otimes B' + A' \otimes B - A' \otimes B'$ obeys
Landau's identity
\begin{equation}
\mathcal{C}^{2} \;=\; 4\,\mathbb{I} \otimes \mathbb{I} \;-\; [A, A'] \otimes [B, B'],
\label{eq:landau}
\end{equation}
whence the quantum bound $\|\mathcal{C}\| \le 2\sqrt{2}$
\autocite{landau1987, tsirelson1980}, against the classical bound $2$
\autocite{chsh1969} and the algebraic maximum $4$. The $2$-adic valuations of the
three rungs are $v_2(2) = 1$, $v_2(2\sqrt{2}) = 3/2$, $v_2(4) = 2$: the quantum bound
occupies the unique half-integral rung, and half-integral valuations are precisely the
signature of ramification --- $\sqrt{2}$ generates the ramified quadratic extension
$\mathbb{Q}_2(\sqrt{2})/\mathbb{Q}_2$. The algebra concurs on where the excess lives.
For qubit observables $A = \mathbf{a} \cdot \boldsymbol{\sigma}$ the commutators lie
in the bivector sector, and the even Clifford algebra of the rational observable form
$\langle 1, 1, 1 \rangle$ is
$C_0(\langle 1,1,1 \rangle) \cong (-1,-1)_{\mathbb{Q}}$ \autocite{lam2005} --- the
ledger's algebra, division at exactly $\{2, \infty\}$. The classical bound consumes
only the commutative data in \eqref{eq:landau}; the excess factor $\sqrt{2}$ is
contributed by the quaternionic sector whose rational form ramifies at two. The
valuation bookkeeping above is exact; the causal reading --- that the
quantum--classical gap is ramification at the binary place made Archimedean-visible
--- is a fenced RHYME, and Problem~\ref{op:padic-tsirelson} states what would cash it.

\begin{openproblem}[The Tsirelson exponent from local invariants]\label{op:padic-tsirelson}
Derive the exponent $3/2$ --- equivalently the excess factor $\sqrt{2}$ --- from the
local invariants of the observable algebra rather than from the operator norm: a
Hilbert-symbol or ramification computation in which the quantum bound's residence in
$\mathbb{Q}_2(\sqrt{2})$ is forced by the fact that the commutator sector generates
$(-1,-1)_{\mathbb{Q}}$. This sharpens the standing question of whether the adelic
product formula constrains the Tsirelson bound from a slogan to a target computation;
any derivation in which $\sqrt{2}$ enters as a ramification datum, however special the
setting, would constitute progress.
\end{openproblem}

\subsection{Loops, solenoids, and the dual of the dyadic tower}
\label{app:padic-solenoid}

The continuity appendix already meets the solenoid, as the witness that connected does
not imply path-connected, excluded at the record layer by Cartan's theorem once the
scalars are real (Remark~\ref{rem:witnesses}). Here is the same object from the loop
side. A closed loop whose traversal count is tallied $p$-adically is the $p$-solenoid
$\Sigma_p = (\mathbb{R} \times \mathbb{Z}_p)/\mathbb{Z}$, the inverse limit of the
circle under $z \mapsto z^{p}$: a compact connected abelian group fibring over the
circle with Cantor fibre $\mathbb{Z}_p$ \autocite{vandantzig1930, hewittross1963}. Its
bonding map is the monomial dynamics at the centre of the $p$-adic dynamical systems
programme \autocite{khrennikovnilsson2004}, and winding $p^{n}$ times is small: many
traversals converge rather than accumulate, which is Proposition~\ref{prop:carry} in
geometric dress. Completing at every prime at once gives the adelic solenoid
$(\mathbb{R} \times \widehat{\mathbb{Z}})/\mathbb{Z}$, the Pontryagin dual of the
discrete rationals \autocite{hewittross1963} --- the character group of exactly the
field whose places the product formula ranges over.

One duality sharpens the imagery into an interface with the continuity appendix. The
phase group is constructed there through the dyadic tower: $U(1)$ is obtained as the
Archimedean closure of the Pr\"ufer group $\mathbb{Z}(2^{\infty})$,
equation~\eqref{eq:dyadic-tower}. Pontryagin duality assigns that same tower a
canonical non-Archimedean completion: the dual of $\mathbb{Z}(2^{\infty})$ is
$\mathbb{Z}_2$ \autocite{hewittross1963}. The very group from which the framework
builds the observer's phase already carries a $2$-adic shadow by duality --- the prime
two's second, independent appearance, after Theorem~\ref{thm:ledger}(ii). Whether this
is coincidence or structure is Problem~\ref{op:padic-solenoid}. The status of this
subsection short of that problem is RHYME, fenced.

\begin{openproblem}[Circle or solenoid]\label{op:padic-solenoid}
The exclusions that force the phase group to be $U(1)$ at the record layer --- no
small subgroups, and Cartan's theorem over the real scalars
(Remark~\ref{rem:witnesses}, \S\ref{app:cont-scalar}) --- consume Archimedean,
record-side premises. Determine whether any pre-record consistency condition of the
framework distinguishes $U(1)$ from its solenoidal completions: is the pre-FANOUT
phase group the circle, the $2$-solenoid, or the adelic solenoid, with the record
layer seeing only the Archimedean face in every case? An answer in either direction is
structural. The circle would mean the Archimedean selection reaches deeper than
records; a solenoid would give the ledger of this appendix a geometric carrier, and
would make the fibrewise digits of the phase literal.
\end{openproblem}

\subsection{Bridge condition, and an invitation}
\label{app:padic-bridge}

One condition governs every correspondence above. Identifying an Archimedean
quasi-probability value with a $p$-adic one requires the value to lie in $\mathbb{Q}$,
or in a fixed number field, so that every completion sees the same element. The
framework's finite, stabiliser-level bookkeeping satisfies this; the Tsirelson datum
satisfies it in $\mathbb{Q}(\sqrt{2})$, whose places above two carry the ramification
reading of \S\ref{app:padic-tsirelson}. Quantities irrational without qualification
must first be located in a number field before the ledger applies; that is a
constraint, and we state it as one.

The six problems are stated so that a specialist can engage each without adopting the
framework whole: Problem~\ref{op:padic-frequency} is a question inside $p$-adic
frequency theory; \ref{op:padic-hensel} inside $p$-adic dynamics;
\ref{op:padic-ghz} inside $p$-adic models of correlation experiments;
\ref{op:padic-reciprocity} inside the arithmetic of quadratic forms;
\ref{op:padic-tsirelson} a finite computation with quaternion algebras; and
\ref{op:padic-solenoid} inside topological group theory, run against the continuity
appendix's premises. The communities they address --- the $p$-adic mathematical
physics of Khrennikov, Volovich, Dragovich and Zelenov
\autocite{khrennikov1995, vladimirovvolovich1989, dragovich1995}, the division-algebra
particle physics of Dixon, Furey, Krasnov and Dubois-Violette, and the $p$-adic
state-space geometry of Palmer's invariant set programme \autocite{palmer2009} ---
have developed with almost no contact, and neither has had reason to ask the other's
questions. Ostrowski's fork, read as complementarity rather than rivalry, is the
meeting point this appendix proposes. What was sterile as a competition --- which
completion is the true home of quantum probability --- becomes a programme as a
division of labour: records at the Archimedean place, ledger at the finite places, and
the product formula as the treaty between them.

% ---------------------------------------------------------------------
% CITATION KEYS (NEW = supplied in refs_padic_ctc_division_2026-07-23T0617.ris;
% EXISTING = assumed in master bibliography, VERIFY each key)
%
% NEW keys:
% serre1973: J.-P. Serre, A Course in Arithmetic, GTM 7, Springer
%   (1973). Hilbert symbols, reciprocity, isotropy in >= 5 variables.
% lam2005: T. Y. Lam, Introduction to Quadratic Forms over Fields,
%   GSM 67, AMS (2005). Levels, u-invariants, even Clifford algebras.
% springerveldkamp2000: T. A. Springer, F. D. Veldkamp, Octonions,
%   Jordan Algebras and Exceptional Groups, Springer (2000).
% jacobson1958: N. Jacobson, Composition algebras and their
%   automorphisms, Rend. Circ. Mat. Palermo 7, 55-80 (1958).
% gouvea1997: F. Q. Gouvea, p-adic Numbers: An Introduction, 2nd ed.,
%   Springer (1997). (Accent in surname carried in the RIS file.)
% artinschreier1927: E. Artin, O. Schreier, Algebraische Konstruktion
%   reeller Koerper, Abh. Math. Sem. Univ. Hamburg 5, 85-99 (1927).
% brauerhassenoether1932: R. Brauer, H. Hasse, E. Noether, Beweis
%   eines Hauptsatzes in der Theorie der Algebren, J. reine angew.
%   Math. 167, 399-404 (1932).
% alberthasse1932: A. A. Albert, H. Hasse, A determination of all
%   normal division algebras over an algebraic number field, Trans.
%   AMS 34, 722-726 (1932).
% landau1987: L. J. Landau, On the violation of Bell's inequality in
%   quantum theory, Phys. Lett. A 120, 54-56 (1987).
% khrennikovnilsson2004: A. Yu. Khrennikov, M. Nilsson, p-adic
%   Deterministic and Random Dynamics, Kluwer (2004).
% vandantzig1930: D. van Dantzig, Ueber topologisch homogene Kontinua,
%   Fund. Math. 15, 102-125 (1930).
% hewittross1963: E. Hewitt, K. A. Ross, Abstract Harmonic Analysis I,
%   Springer (1963). Solenoids, Pruefer duals, character groups.
% palmer2009: T. N. Palmer, The invariant set postulate, Proc. R.
%   Soc. A 465, 3165-3185 (2009).
% tsirelson1980: B. S. Cirel'son, Quantum generalizations of Bell's
%   inequality, Lett. Math. Phys. 4, 93-100 (1980). (Included in the
%   RIS batch; VERIFY against master bibliography for duplication.)
% chsh1969: J. F. Clauser, M. A. Horne, A. Shimony, R. A. Holt,
%   Proposed experiment to test local hidden-variable theories, PRL
%   23, 880-884 (1969). (Included in RIS; VERIFY for duplication.)
%
% EXISTING keys (VERIFY spelling of each against master bibliography):
% ostrowski1916: as in app:continuity-open.
% thomson1954: as in app:continuity-open.
% khrennikov1995: A. Yu. Khrennikov, p-adic description of Dirac's
%   hypothetical world with negative probabilities, IJTP 34 (1995);
%   listed in the adelic open-problem document.
% khrennikov2009: A. Yu. Khrennikov, p-adic probability prediction of
%   correlations..., arXiv:0906.0509; listed in the adelic document.
% vladimirovvolovich1989, dragovich1995: adelic QM; listed in the
%   adelic document.
% hurwitz1898, adams1960: listed in the adelic document and main text.
% deutsch1991, aaronsonwatrous2009, laraudogoitia1998, spekkens2008:
%   listed in the grounding-argument reference list.
% =====================================================================
